\scanpage{61}
% block: reconstruction continued
$\handwritten{\mathcal{M}}$ en termes de la \unsure{famille|faamille} des $A_i$ \deleted{et des entiers $d_i=\operatorname{rang}E_i$} comme suit\deleted{e}. On considère d'abord la catégories sous-additive \deleted{engendrée} ayant comme objets les éléments $\handwritten{\underline{\unsure{\gamma|Y}}}$ de $\underline{N}^{(I)}$, ($\handwritten{\underline{\gamma}}=(\handwritten{\gamma}_i)_{i\handwritten{\in}I}$ étant interprété comme $\handwritten{\bigotimes}E_i^{\handwritten{\gamma}_i}$), \deleted{ou} avec $\operatorname{Hom}(\handwritten{\underline{\gamma}},\handwritten{\underline{\gamma}'})=0$ si $\handwritten{\underline{\gamma}\ne\underline{\gamma}'}$ et $=\bigotimes_{\handwritten{i}} A_i^{\otimes\handwritten{\gamma_i}}$ si $\handwritten{\underline{\gamma}=\underline{\gamma}'}$, la composition des morphismes évidentes. On définit de façon évidente le foncteur $\otimes$ (addition de $\underline{N}^{(I)}$ sur les objets, produit tensoriel sur les homomorphismes) et ses contraintes,\margin{\handwritten{\unsure{attention à} contrainte de commutativité}} de façon que l'on trouve une sous-$\otimes$-catégorie. Son enveloppe Karoubienne est une\inserted{$\otimes$-}catégorie abélienne, qui s'identifie à la sous-catégorie pleine de $\handwritten{\mathcal{M}}$ formée des éléments dont les composantes homogènes non nulles sont à degrés $\handwritten{\in}\underline{N}^{(I)}$,\inserted{et a un foncteur fibre sur \deleted{$k$}\inserted{$\overset{\handwritten{-}}{k}$.}} D'autre part, pour tout $i$, si $A_i$ est de \deleted{degré}\inserted{rang} $d_i^2$ (donc $E_i$ de rang $d_i$), soit $L_i=\det E_i={\handwritten{\unsure{A|\bigwedge}_{\unsure{2|k}}}}^{d_i}\handwritten{E_i}$ (où $E_i$ désigne l'objet de la catégorie construite correspondant \deleted{à \unsure{il|le}} $i$.ème élément de la base canonique de $\underline{Z}^{(I)}$, soit $1_i$ ...), c'est un objet préinversible ($\inserted{\handwritten{X\to}}\deleted{\unsure{\otimes}}L_i\otimes\handwritten{X}$ est pleinement fidèle) et en inversant les $L_i$, on trouve (à \deleted{\unsure{équ}}équivalence près\deleted{)} respectant tout) la catégorie $\handwritten{\mathcal{M}}$.

% block: nonfree case
Quand $\handwritten{\unsure{\Gamma|T}}$ n'est pas libre, on trouve une construction analogue en termes d'un système de générateurs\inserted{convenable} $(\handwritten{\gamma}_i)_{i\handwritten{\in}I}$, mais en``\unsure{renant|rendant} \deleted{de plus} égaux à 1 \deleted{x}'' de plus certains éléments.\margin{A expliciter !} Ces constructions sont bien adaptées pour décrire les $\otimes$-foncteurs de $\handwritten{\mathcal{M}}$ dans une \deleted{$\otimes$-}catégorie tannakienne, et pour décrire le changement de base.( Pour cette dernière raison notamment, il aurait été maladroit dans le cas général de prendre les $M_i$ simples ....). Attention, dans la description des $\otimes$-foncteurs en question, de tenir compte des contraintes, notamment de celle de commutativité; il semble qu'il doive revenir au même d'exiger que le rang est conservé ...
% archive sheet number: 60
